% ==============================================================================
% Discussion
% ==============================================================================

\section{Discussion}
\label{sec:discussion}

\textbf{[Placeholder: to be completed after results are available.  The
following is the planned structure.]}

\paragraph{The deliberation gradient.}
[Discuss: do the results support a graded deliberation-intensity dimension
rather than a binary S1/S2 dichotomy?  Is the gradient continuous across
layers?  Does it vary smoothly with task difficulty or show sharp transitions?
Connect to \citet{melnikoff2018mythical}'s critique of the binary framework
and \citet{ziabari2025reasoning}'s ``Reasoning on a Spectrum'' finding.]

\paragraph{Standard vs.\ reasoning-trained models.}
[Discuss: how does reasoning distillation change the deliberation gradient?
Does R1-Distill-Llama-8B show (a) a stronger S2 signal at the same layers as
Llama-3.1-8B-Instruct, (b) a qualitatively different layer-wise profile, or
(c) no meaningful difference?  Relate to the CKA analysis and the
\citet{codaforno2025dual} finding of overlapping subspaces.]

\paragraph{Comparison with CogBias.}
[Discuss: how do our probe results compare with \citet{huang2026cogbias}?
Our probes measure S1/S2 processing mode; theirs target bias susceptibility for
mitigation.  Do the two approaches identify the same layers and directions?
Complementary contributions: they demonstrate practical bias reduction via
steering; we provide a richer mechanistic characterization of the underlying
computation.]

\paragraph{Implications for AI safety.}
[Discuss: understanding when models engage in heuristic versus deliberative
processing has direct implications for trust calibration.  If we can detect
from internal representations whether a model is ``actually reasoning'' on a
given input (vs.\ pattern-matching), this enables more targeted oversight.
However, note the limitation that chain-of-thought faithfulness is low
\citep{jiang2025aha}---the activation-level signal may be more trustworthy than
the token-level signal, but we cannot yet guarantee that high deliberation
intensity entails faithful reasoning.]

\paragraph{Limitations.}
[Discuss: (i)~Model scale: our largest model is 9B parameters.  The
deliberation gradient may manifest differently in models with 70B+ parameters.
(ii)~Benchmark scope: 7 task categories and 284 items cover a limited range of
cognitive biases.  (iii)~SAE coverage: pre-trained SAEs may not capture all
relevant features, especially for the distilled models.  (iv)~Causal evidence:
activation steering demonstrates causal relevance of directions, not circuits.
(v)~Anthropomorphism risk: our ``deliberation intensity'' language is
operational, not a claim about phenomenology.]

\paragraph{Future work.}
[Discuss: (i)~Scaling: apply the framework to 70B+ models.
(ii)~Broader tasks: extend beyond cognitive biases to other domains where
heuristic/deliberative processing differs (e.g., moral reasoning, scientific
reasoning).  (iii)~Circuit-level analysis: move from feature-level to
circuit-level understanding of the deliberation gradient.  (iv)~Temporal
dynamics: finer-grained analysis of how deliberation intensity evolves across
the thinking trace.  (v)~Cross-species comparison: compare the deliberation
gradient in LLMs to neural signatures of dual-process cognition in humans
\citep{deneys2023dual}.]
